
\section{Charge extortion \ddc}
\label{chap5:sect:chargeExtortion}
This section explains the charge extortion mechanism at work during BBI which allows fault creation.
The voltage pulse generator, at each edge of its pulse, injects and then extorts electrical charges into and out of the IC.
%\subsection{Logic considerations \ddc}
%\label{chap5:sect:chargeExtortion:subsect:logicConsiderations}
%
%\begin{figure}[htbp!]
    \centering
    \includegraphics[width=0.6\textwidth]{5_faultModel/figures/dffLogic.pdf}
    \caption{D Flip-Flop logic schematic}
    \label{chap:5;fig:dffLogic}
\end{figure}
%In order to understand electrically why faults occur, it is important to begin with logic considerations.
%Nowadays, the vast majority of commercial ICs are sequential.
%The main building block of a sequential IC is the Edge-Triggered D Flip-Flop (DFF).
%Fig. \ref{chap:5;fig:dffLogic} presents the logic schematic of a DFF.
%Because DFFs work on clock edges, it is interesting to target them at specific times when performing fault injection.
%\textcolor{magenta}{Paragraphe à terminer, compléTer, revoir...}
%%The block has two inputs and two outputs.
%%This building block is ruled by a clock, input at C.
%%At each rising edge of the clock, the output samples the input value.
%%It is thus a memory cell, which stores at Q for a short amount of time the value present at D.
%%What is important to remember here is that outside of the rising edge, any change at D is not seen at Q.
%%Therefore, when performing fault injections, it is interesting to affect the value at D long enough to let the DFF sample an incorrect value.
%This fault model has been proposed in \cite{emfiFF} for EMFI.
%
%\textcolor{magenta}{Parler du fait que les circuits séquentiels échantillonnent périodiquement les valeurs des DFFs (registres). Du coup, si on change un niveau logique suffisamment longtemps pour qu'il perdure pendant l'échantillonnage, c'est gagné ! Faute !, sinon Non.}

\subsection{Sequential logic operation and simple fault model \ddc}
\label{chap5:sect:chargeExtortion:subsect:seqLogic}
As sequential logic is ubiquitous in contemporary integrated circuits, we shall examine its fundamental workings in greater detail.
Sequential logic relies on a core element: the edge-triggered D flip-flop (DFF).
They are a memory component that is governed by a clock.
At each rising-edge or falling-edge (depending on the design) of the clock, DFFs sample their input and replicate it at their output.
Between DFFs are placed the logic gates, which fulfill a specific logical function.
Because of this, values in sequential logic circuits can only be changed at the clock edges.
Hence, in the event that an adversary is able to alter a logical value for an extended period of time at the input of a DFF, the subsequent combinatorial logic will yield an incorrect value, which will propagate to the subsequent DFFs.

\begin{figure}[H]
    \centering
    \includegraphics[width=15cm]{5_faultModel/figures/bbiFaultSusceptibility.pdf}
    \caption{BBI sampling fault susceptibility}
    \label{chap:5;fig:bbiSusc}
\end{figure}
This fault model we described previously is depicted in Fig. \ref{chap:5;fig:bbiSusc} and was first introduced in \cite{emfiFF} for EMFI.

\subsection{Charge extortion \ddc}
\label{chap5:sect:chargeExtortion:subsect:chargeExtortion}
